% SLICESPREADTRANSFER1 2026-10-17 20:45 Asia/Seoul
\documentclass[11pt]{article}
\usepackage[T1]{fontenc}
\usepackage{lmodern,amsmath,amssymb,amsthm,booktabs,hyperref}
\usepackage[margin=0.9in]{geometry}
\newtheorem{theorem}{Theorem}
\newtheorem{lemma}[theorem]{Lemma}
\newtheorem{corollary}[theorem]{Corollary}
\newtheorem{definition}[theorem]{Definition}
\newcommand{\EXACT}{\textbf{PROVED/EXACT}}
\newcommand{\KNOWN}{\textbf{KNOWN THEOREM}}
\newcommand{\OPEN}{\textbf{OPEN}}
\newcommand{\RET}{\textbf{FALSIFIED/RETIRED}}
\title{SLICESPREADTRANSFER1: Spread-DNF Transfer Without a Term-Count Loss}
\author{P versus NP project}
\date{17 October 2026}
\begin{document}\maketitle

\noindent\textbf{Status.} The exact-slice spread-DNF second-moment theorem, the relative-moment transfer bound, and the spread-cluster corollary are \EXACT at their stated models, pending independent review. Rao's spread-matching lower-bound method, Cavalar--Kumar--Rossman's closure/robust-sunflower approximation, and Kahn--Linial--Samorodnitsky/Sherstov approximate inclusion--exclusion are \KNOWN controls. A circuit-to-bounded-width low-spread-cluster approximation theorem remains \OPEN. Consequently $P\ne NP$ remains \OPEN.

\section{IDEAS PROPOSED BEFORE LITERATURE}
\paragraph{A: remove the DNF term-count loss.} The trivial implication from moment matching to DNF fooling loses a factor equal to the number of terms. Search for an exact bound depending instead on a spread parameter of the term family. Falsifier: a spread family with small negative acceptance but large first moment.
\paragraph{B: spread-cluster decomposition.} If a term family is not globally spread, factor dense cores and ask whether the number of resulting spread clusters is a better mass-production parameter than the number of terms. Missing theorem: small monotone circuits should admit an approximation with small spread-cluster number.
\paragraph{C: arbitrary-polynomial-exponent audit.} The NAEGAUGEHIGHGIRTH1 positive distribution only matches moments through $c_*\log q$, $c_*\approx0.2094$. Test whether a spread transfer can make this logarithmic window sufficient for every fixed circuit exponent $K$ without requiring width proportional to $K\log q$.
\paragraph{D: unsigned NAE as a hostile alternate target.} Positive NAE-3-SAT-E4 already lives on an unsigned exact hypergraph layer. This removes the fixed-support sign-parity shortcut but also removes the gauge-code pseudorandomness theorem. It is retained as a hostile comparison, not substituted into the current proof spine in this run.

\section{LITERATURE CHECK AND REVISIONS}
Rao's 2026 spread-matching method shows that tailored spread structure can yield strong monotone circuit lower bounds that are invisible to generic robust-sunflower counting. The theorem is matching-specific and uses a structured negative distribution; it does not directly transfer to the SAT/NAE slice.

Cavalar--Kumar--Rossman's approximation framework compresses small monotone circuits by closure and robust-sunflower trimming. In its standard gate-by-gate form the accumulated error is bounded by circuit size times an exponentially small width-dependent error, so handling a size-$N^K$ circuit in that template typically requires the width constant to grow with $K$. This clashes quantitatively with the fixed coefficient $c_*$ in the presently proved NAE moment window.

Kahn--Linial--Samorodnitsky and later Sherstov show that low-order inclusion--exclusion information need not approximate an arbitrary union unless the order is much larger than logarithmic in the number of events in the worst case. Therefore low moments alone cannot fool arbitrary polynomial-size DNFs; a structural hypothesis such as spread is essential.

These controls motivate Candidate A rather than another generic moment argument.

\section{CONSTRUCTION OR PROOF ATTEMPT}
\subsection{Spread DNF on the uniform exact slice}
Let $\nu$ be the uniform distribution on the exact layer $\binom{[N]}t$, put
\[p=t/N,\]
and let $\mathcal F$ be a family of $L$ distinct $w$-sets. Define the monotone DNF
\[D_{\mathcal F}(x)=\bigvee_{T\in\mathcal F}\prod_{i\in T}x_i,\qquad R(x)=\sum_{T\in\mathcal F}\prod_{i\in T}x_i.\]

\begin{definition}[$\rho$-spread term family]
$\mathcal F$ is $\rho$-spread if for every $A\subseteq[N]$,
\[d_{\mathcal F}(A):=|\{T\in\mathcal F:A\subseteq T\}|\le L\rho^{-|A|}.\]
Put $\theta=\rho p$ and $\lambda=Lp^w$.
\end{definition}

\begin{theorem}[Exact-slice spread second moment]
Assume $w\le t/2$. Set
\[c_0=\exp\!\left(-\frac{w(w-1)}t\right).\]
Then
\[\mathbb E_\nu R\ge c_0\lambda,\qquad \mathbb E_\nu(R)_2\le \lambda^2\exp(w/\theta),\]
and therefore
\[\boxed{\Pr_\nu[D_{\mathcal F}=1]\ge \frac{c_0^2\lambda}{1+e^{w/\theta}\lambda}.}\]
\end{theorem}
\begin{proof}
For a fixed $w$-set,
\[h_w:=\Pr_\nu[T\subseteq X]=\frac{(t)_w}{(N)_w}.\]
Since $w\le t/2$,
\[\frac{h_w}{p^w}=\prod_{i=0}^{w-1}\frac{1-i/t}{1-i/N}\ge\prod_{i=0}^{w-1}(1-i/t)\ge e^{-w(w-1)/t}=c_0,\]
using $\log(1-x)\ge-2x$ for $x\le1/2$. Thus $\mathbb E R=Lh_w\ge c_0\lambda$.

For distinct terms $T,S$, the exact-slice containment probability is at most $p^{|T\cup S|}$. For fixed $T$,
\[\sum_{S\in\mathcal F}p^{|S\setminus T|}\le\sum_{A\subseteq T}d_{\mathcal F}(A)p^{w-|A|}\le Lp^w\left(1+\frac1{\rho p}\right)^w\le\lambda e^{w/\theta}.\]
Hence summing over $T$ gives $\mathbb E(R)_2\le\lambda^2e^{w/\theta}$. Also $\mathbb E R\le\lambda$. Paley--Zygmund at threshold zero yields
\[\Pr[R>0]\ge\frac{(\mathbb E R)^2}{\mathbb E R^2}\ge\frac{c_0^2\lambda^2}{\lambda+e^{w/\theta}\lambda^2},\]
as claimed.
\end{proof}

\subsection{Relative-moment transfer without the factor $L$}
Let $\mu$ be another distribution on $\{0,1\}^N$ satisfying the upper relative moment condition
\[\Pr_\mu[T\subseteq X]\le(1+\delta)\Pr_\nu[T\subseteq X]\]
for every $T\in\mathcal F$.

\begin{theorem}[Spread-DNF transfer]
Let $\eta=\Pr_\nu[D_{\mathcal F}=1]$. If
\[c_0^2>\eta e^{w/\theta},\]
then
\[\boxed{\Pr_\mu[D_{\mathcal F}=1]\le (1+\delta)\frac{\eta}{c_0^2-\eta e^{w/\theta}}.}\]
In particular, if $w^2=o(t)$, $w/\theta=O(1)$, $\delta=o(1)$, and $\eta=o(1)$, then
\[\Pr_\mu[D_{\mathcal F}=1]=O(\eta),\]
with no dependence on the number $L$ of DNF terms.
\end{theorem}
\begin{proof}
The preceding lower bound gives
\[\eta(1+e^{w/\theta}\lambda)\ge c_0^2\lambda,\]
hence
\[\lambda\le\frac{\eta}{c_0^2-\eta e^{w/\theta}}.\]
On the other hand, a union bound under $\mu$ gives
\[\Pr_\mu[D=1]\le\sum_{T\in\mathcal F}\Pr_\mu[T]\le(1+\delta)Lh_w\le(1+\delta)\lambda.\]
Combine the two inequalities.
\end{proof}

This is the main advance of the run. The earlier trivial moment-transfer estimate was $(1+\delta)L\eta$. Spread removes the entire term-count factor.

\subsection{Immediate application to NAEGAUGEHIGHGIRTH1}
Let $\mu_q^+$ be the exact-density-$4/3$ UNSAT NAE distribution constructed in NAEGAUGEHIGHGIRTH1 and let $\nu_q$ be the uniform exact-density layer. For every $c<c_*$, $c_*\approx0.20940225$, the project proved uniformly for all coordinate sets $T$ with $|T|\le c\log q$ that
\[\Pr_{\mu_q^+}[T]=(1+o(1))\Pr_{\nu_q}[T].\]

\begin{corollary}[Spread-DNF invisibility in the proven moment window]
Fix constants $c<c_*$ and $A>0$. Let $\mathcal F_q$ be a sequence of $\rho_q$-spread term families of width
\[w_q\le c\log q,\]
such that
\[\rho_q\frac{t}{N}\ge A w_q.\]
If
\[\Pr_{\nu_q}[D_{\mathcal F_q}=1]=o(1),\]
then
\[\boxed{\Pr_{\mu_q^+}[D_{\mathcal F_q}=1]=o(1).}\]
\end{corollary}
\begin{proof}
Here $t=\Theta(q)$ and $w_q=O(\log q)$, so $w_q^2/t=o(1)$ and $c_0=1-o(1)$. The spread assumption gives $w_q/\theta_q\le1/A$. Moment matching supplies $\delta=o(1)$. Apply the transfer theorem.
\end{proof}

This corollary can handle arbitrarily many DNF terms. Hence the finite coefficient $c_*$ no longer creates an immediate barrier from a polynomial number of terms.

\subsection{Spread-cluster number}
Global spread is too strong for arbitrary term families because many terms may share a dense core. Define a cluster to be a pair $(C,\mathcal G)$ representing terms $C\cup R$ with $R\in\mathcal G$, where the residual family $\mathcal G$ is $\rho$-spread. Let
\[\kappa_\rho(\mathcal F)\]
be the minimum number of disjoint term clusters whose union is $\mathcal F$.

\begin{lemma}[Existence of a spread-cluster decomposition]
Every finite $w$-uniform family admits such a decomposition with at most $|\mathcal F|$ clusters.
\end{lemma}
\begin{proof}
Start with empty core. If the residual family is $\rho$-spread, stop. Otherwise choose a nonempty set $A$ violating the spread inequality, branch off all residual terms containing $A$, add $A$ to that branch's core, and recurse both on that branch after stripping $A$ and on the remaining terms. Core size strictly grows along every nonterminal branch; the process terminates. Distinct leaves contain disjoint sets of original terms, so there are at most $|\mathcal F|$ leaves.
\end{proof}

\begin{corollary}[Cluster-transfer bound, schematic exact constants suppressed]
Suppose relative moment matching holds for every full term and every cluster core through width $w$. If every residual cluster has $w'/\theta=O(1)$ and $w'^2=o(t)$, then for negative DNF acceptance $\eta=o(1)$, each cluster contributes $O(\eta)$ positive acceptance. Consequently
\[\Pr_\mu[D_{\mathcal F}=1]\le O(\kappa_\rho(\mathcal F)\eta).\]
\end{corollary}
\begin{proof}
For a cluster core $C$, write $a=\Pr_\nu[C]$. If $a=O(\eta)$, then the positive cluster is at most $\Pr_\mu[C]=(1+o(1))a=O(\eta)$. Otherwise condition on $C$. The conditional negative cluster probability is $\eta_C\le\eta/a=o(1)$; the residual family is spread, so the spread-DNF transfer argument bounds the sum of its positive full-term probabilities by $O(\eta_C)$. Multiplying by the core probability gives $O(\eta)$. Sum over the disjoint term clusters.
\end{proof}

The important new structural parameter is therefore not DNF term count but spread-cluster number. A useful circuit approximation theorem may allow enormous mass-produced term families as long as their cluster number is small enough relative to the negative error.

\subsection{Arbitrary polynomial exponent audit}
Suppose a size-$N^K$ monotone circuit can be transformed, for every desired constant $B$, into a bounded-width DNF $D$ satisfying
\[\Pr_\nu[C\ne D]\le N^{-B},\qquad \Pr_{\mu^+}[C\ne D]\le N^{-B},\]
with width $O_K(1)$ and
\[\kappa_\rho(D)\le N^{b_K}\]
for some finite $b_K$, while every residual cluster has $w'/\theta=O_K(1)$. Choosing $B>b_K$ makes the cluster-transfer error $o(1)$. Since the NAE moment window grows like $c_*\log q$, every fixed width $O_K(1)$ eventually lies inside it.

Therefore the arbitrary-exponent issue no longer requires one universal numerical parameter tuple. It reduces to the qualitative theorem:
\[\boxed{\text{poly-size monotone circuit}\Longrightarrow\text{constant-width, polynomial-spread-cluster approximation with arbitrarily small polynomial error}.}\]
This theorem is \OPEN. Standard gate-by-gate CKR trimming does not currently provide it at the required same-layer positive error, but Rao-style spread decompositions suggest the correct structural language.

\section{TESTS AND OUTCOME}
Three hostile controls clarify the exact scope.
\begin{enumerate}
\item A complete family of all $w$-sets can be highly spread, but its negative acceptance is one; the denominator condition of the transfer theorem correctly gives no useful conclusion.
\item A large star of terms sharing one core is not globally spread. The cluster formulation isolates the core instead of falsely applying the global theorem.
\item Kahn--Linial--Samorodnitsky/Sherstov examples show that low moments do not fool arbitrary large unions, so the spread or low-cluster hypothesis cannot be dropped.
\end{enumerate}

Outcome: the project now has a model-matched theorem converting the proved logarithmic NAE moment matching into invisibility for a broad class of bounded-width DNFs with no dependence on their number of terms. This removes one major quantitative obstruction. The remaining load-bearing task is a deep-circuit structural theorem controlling spread-cluster number after approximation. Raw minterm count, cancellation order, additive HD mass, and fixed-support parity restrictions remain retired.

\paragraph{Next target.}
Audit CKR closure/trimming and Rao's spread-matching decomposition specifically for spread-cluster number. Determine whether a gate-local recurrence for $\log\kappa_\rho$ or a global decomposition theorem survives union/intersection mass production. In parallel, use the signed-orbit DAG machinery only to compute extremal cluster cores and codegrees for bounded-width approximators; do not infer hardness from DAG state count.

Root \texttt{main.tex} remains truncated and was not modified.

\begin{thebibliography}{5}
\bibitem{Rao26} A. Rao, \emph{Monotone Circuit Lower Bounds from Spread Matchings}, ECCC TR26-129, 2026.
\bibitem{CKR22} B. P. Cavalar, M. Kumar, B. Rossman, \emph{Monotone Circuit Lower Bounds from Robust Sunflowers}, Algorithmica 84, 2022.
\bibitem{KLS96} J. Kahn, N. Linial, A. Samorodnitsky, \emph{Inclusion-exclusion: Exact and approximate}, Combinatorica 16, 1996.
\bibitem{Sherstov18} A. Sherstov, \emph{The power of asymmetry in constant-depth circuits}, related sharp approximate-inclusion--exclusion refinements; used here only as a hostile-control literature pointer.
\bibitem{ProjectNAE} Project report NAEGAUGEHIGHGIRTH1, 17 October 2026.
\end{thebibliography}
\end{document}
% FINAL_VERIFICATION_SENTINEL::MAIN::SLICESPREADTRANSFER1::2026-10-17_2045::P_VS_NP_OPEN